\section{几何向量的概念及其线性运算}

\section{几何向量的数量积、向量积和混合积}

\subsection{向量在轴上的投影}

\subsection{几何向量的数量积}
\begin{definition}
设$\boldsymbol{a},\boldsymbol{b}$是两个几何向量，称$|\boldsymbol{a}|\,|\boldsymbol{b}|\cos\theta\;(\theta=\langle\boldsymbol{a},\boldsymbol{b}\rangle)$为$\boldsymbol{a}$与$\boldsymbol{b}$的**数量积**或**内积**，记作$\boldsymbol{a}\cdot\boldsymbol{b}$。
\end{definition}

\subsection{几何向量的向量积}
对空间中任意两个向量$\boldsymbol{a}$和$\boldsymbol{b}$，我们来确定第三个向量，叫做$\boldsymbol{a}$与$\boldsymbol{b}$的向量积，用$\boldsymbol{a}\times\boldsymbol{b}$表示（注意次序），其定义如下：

\begin{definition}
若$\boldsymbol{a},\boldsymbol{b}$不平行，则
\begin{enumerate}
    \item $|\boldsymbol{a}\times\boldsymbol{b}|=|\boldsymbol{a}|\,|\boldsymbol{b}|\sin\theta,\quad \theta=\langle\boldsymbol{a},\boldsymbol{b}\rangle$；
    \item %待补
\end{enumerate}
\end{definition}

几何向量的向量积具有下列性质：
\begin{enumerate}
    \item $\boldsymbol{a}\times\boldsymbol{b}=-\boldsymbol{b}\times\boldsymbol{a}$；
    \item $\boldsymbol{a}\times\boldsymbol{b}\perp \boldsymbol{a},\ \boldsymbol{a}\times\boldsymbol{b}\perp \boldsymbol{b}$；
    \item %待补
\end{enumerate}

注意：%待补

\subsection{几何向量的混合积}
\begin{definition}
已知三个向量$\boldsymbol{a},\boldsymbol{b}$和$\boldsymbol{c}$，先作$\boldsymbol{a}$和$\boldsymbol{b}$的向量积$\boldsymbol{a}\times\boldsymbol{b}$，把所得的向量与$\boldsymbol{c}$再作数量积$(\boldsymbol{a}\times\boldsymbol{b})\cdot\boldsymbol{c}$，这样得到的数量叫做三向量$\boldsymbol{a},\boldsymbol{b},\boldsymbol{c}$的混合积，记为$[\boldsymbol{a}\boldsymbol{b}\boldsymbol{c}]$。
\end{definition}

向量的混合积有如下几何意义：$[\boldsymbol{a}\boldsymbol{b}\boldsymbol{c}]=(\boldsymbol{a}\times\boldsymbol{b})\cdot\boldsymbol{c}$的绝对值表示以向量$\boldsymbol{a},\boldsymbol{b},\boldsymbol{c}$为棱的平行六面体的体积。

\subsection{几何向量的坐标}

\subsection{几何向量的坐标运算}
\begin{enumerate}
    \item[(5)] $\boldsymbol{a}\times\boldsymbol{b}=$ %待补
    \item[(6)] $[\boldsymbol{a}\boldsymbol{b}\boldsymbol{c}]=
\begin{vmatrix}
a_x & a_y & a_z\\
b_x & b_y & b_z\\
c_x & c_y & c_z
\end{vmatrix}$
\end{enumerate}


\section{空间中的平面与直线}

\subsection{空间中平面的方程}
设平面$\pi$通过点$M_0(x_0,y_0,z_0)$并且垂直于非零向量$\boldsymbol{n}=(A,B,C)$。

称垂直于平面$\pi$的非零向量$\boldsymbol{n}=(A,B,C)$为平面$\pi$的**法向量**。

设$M(x,y,z)$是平面$\pi$上的任意一点，则$\overrightarrow{M_0M}$与$\boldsymbol{n}$垂直，从而
\[
A(x-x_0)+B(y-y_0)+C(z-z_0)=0,
\]
上式就是平面$\pi$的方程，称其为平面$\pi$的**点法式方程**。

\subsection{空间中直线的方程}
已知点$M_0(x_0,y_0,z_0)$，并且%待补。

标准方程

参数方程

\subsection{位置关系}
\subsubsection{平面与平面的位置关系}
设有两个平面
\[
\begin{cases}
\pi_1:A_1x+B_1y+C_1z+D_1=0,\\
\pi_2:A_2x+B_2y+C_2z+D_2=0,
\end{cases}
\]
其中$A_i,B_i,C_i$不同时为零，$i=1,2$。
\begin{enumerate}
    \item $\pi_1,\pi_2$重合 %待补
\end{enumerate}

\subsubsection{直线与直线的位置关系}
设有两条直线
\begin{enumerate}
    \item %待补
    \item %待补
    \item $L_1,L_2$交于一点的充要条件是混合积$[\boldsymbol{s}_1\boldsymbol{s}_2\overrightarrow{M_1M_2}]=0$且$\boldsymbol{s}_1,\boldsymbol{s}_2$不平行。
    \item %待补
\end{enumerate}

\subsubsection{直线与平面的位置关系}

\subsection{距离}
\paragraph{点到平面的距离}

\paragraph{平面到平面的距离}
两平行%待补

\paragraph{点到直线的距离}

\paragraph{两平行直线间的距离}

\paragraph{两条异面直线间的距离}
设有两条异面直线：
\[
L_1:\frac{x-x_1}{m_1}%待补
\]

$\overrightarrow{P_1P_2}$在$\boldsymbol{s}_1\times\boldsymbol{s}_2$上的投影的绝对值就是$L_1$与$L_2$之间的距离$d$，即
\[
d=%待补
\]

\subsection{平面束}
称通过给定直线$L$的所有平面的全体为通过直线$L$的**平面束**。

设$L$的一般方程为
\[
A_1x %待补
\]